\chapter{Big sets}

本章研究 set 的大小。有限集合可以数出元素个数；无限集合则需要用
function 比较大小。核心语言是 bijection, injection, cardinality,
Cantor--Bernstein theorem, Cantor's theorem, Axiom of Choice, Zorn's
Lemma, density, and well-ordering。写证明时要分清两件事：一个 set 是否
可以 enumerate，是 cardinality 的问题；一个 subset 是否在 real line 中到处
都能逼近，是 order and metric structure 的问题。

\section{Cardinality}

\begin{definition}[same cardinality]
设 \(X,Y\) 是 sets。若存在 bijection \(f:X\to Y\)，就说 \(X\) 与
\(Y\) have the same cardinality，记作
\[
  \card{X}=\card{Y}.
\]
这里的 bijection 同时要求 injective and surjective：每个 \(x\in X\)
对应唯一的 \(Y\)-element，并且每个 \(Y\)-element 都被某个 \(x\) 命中。
\end{definition}

\begin{definition}[finite, countable, uncountable]
设 \(X\) 是 set。若存在 \(n\in\N\) 使得 \(\card{X}=\card{n}\)，则
\(X\) is finite；此时通常写 \(\card{X}=n\)。若 \(X\) is finite 或
\(\card{X}=\card{\N}\)，则 \(X\) is countable。不是 countable 的 set
称为 uncountable。
\end{definition}

\begin{reviewbox}{How to prove equal cardinality}
证明 \(\card{X}=\card{Y}\) 有两条常用路线。第一条是直接构造
bijection \(X\to Y\)，然后分别验证 injectivity and surjectivity。第二条
是分别构造 injections \(X\to Y\) and \(Y\to X\)，再引用
Cantor--Bernstein theorem。第二条在处理 intervals 和删去少数点的
uncountable sets 时尤其方便。
\end{reviewbox}

\subsection{Countable examples}

\begin{exerciseblock}[Lecture Notes Exercise 62]
Prove that \(\card{\N}=\card{\Z}\).
\end{exerciseblock}

\begin{solution}
本讲义采用 \(\N=\{0,1,2,\ldots\}\)。定义 \(f:\N\to\Z\) 为
\[
  f(0)=0,\qquad
  f(2k+1)=k+1,\qquad
  f(2k+2)=-(k+1)
  \quad (k\in\N).
\]
这个 function 将 natural numbers 依次送到
\[
  0,1,-1,2,-2,3,-3,\ldots .
\]
先证明 \(f\) is injective。若 \(f(n)=0\)，由定义只能有 \(n=0\)。
若 \(f(n)>0\)，则 \(n\) 必为 \(2k+1\)，且 \(f(n)=k+1\) 决定了
\(k\)，从而决定了 \(n\)。若 \(f(n)<0\)，则 \(n\) 必为 \(2k+2\)，且
\(f(n)=-(k+1)\) 也决定了 \(k\)，从而决定了 \(n\)。正整数、负整数和
\(0\) 三类 image 互不重叠，所以不同的 input 不会有相同 image。

再证明 \(f\) is surjective。取任意 \(z\in\Z\)。若 \(z=0\)，则
\(z=f(0)\)。若 \(z=m>0\)，其中 \(m\in\N\) 且 \(m\ge 1\)，则
\[
  z=m=f(2m-1).
\]
若 \(z=-m<0\)，其中 \(m\ge 1\)，则
\[
  z=-m=f(2m).
\]
每个 integer 都出现在 \(f\) 的 image 中。因此 \(f\) is bijective，
所以 \(\card{\N}=\card{\Z}\)。
\end{solution}

\begin{explanation}
这个 proof 的重点不是说两个 infinite sets ``一样多''，而是给出一条
明确的 enumeration。正整数和负整数若分别排在两段，负整数永远等不到；
交替列出 \(1,-1,2,-2,\ldots\) 可以保证每个 integer 在有限步内出现。
\end{explanation}

\begin{exerciseblock}[Worksheet 8 Exercise 1]
Prove that \(\card{\Z}=\card{\N}\).
\end{exerciseblock}

\begin{solution}
由 Lecture Notes Exercise 62 已构造 bijection \(f:\N\to\Z\)。bijection
有 inverse function \(f^{-1}:\Z\to\N\)，而 inverse of a bijection 仍是
bijection。因此 \(\card{\Z}=\card{\N}\)。

若要写出 explicit formula，可以定义
\[
  g(z)=
  \begin{cases}
    0,&z=0,\\
    2z-1,&z>0,\\
    -2z,&z<0.
  \end{cases}
\]
由 Lecture Notes Exercise 62 中的 \(f\) 可见 \(g=f^{-1}\)，所以
\(g\) is bijective。
\end{solution}

\begin{explanation}
Worksheet 8 Exercise 1 与 Lecture Notes Exercise 62 是同一个 cardinality
statement，只是等号两边顺序相反。cardinality equality 是 symmetric：
若 \(X\) 与 \(Y\) 之间有 bijection，则把箭头反过来得到 \(Y\) 与 \(X\)
之间的 bijection。
\end{explanation}

\begin{theorem}[\(\Q\) is countable]
The rational numbers form a countable set.
\end{theorem}

\begin{proof}
先处理 positive rationals。把每个 positive rational 写成 lowest terms
\(p/q\)，其中 \(p,q\in\N\) 且 \(p,q\ge 1\)。将 pair \((p,q)\) 放在
infinite grid 的第 \(p\) 行第 \(q\) 列，然后按 diagonal
\[
  p+q=2,\quad p+q=3,\quad p+q=4,\ldots
\]
依次扫描。每条 diagonal 上只有有限多个 pairs；每个 pair \((p,q)\)
属于唯一的 diagonal \(p+q\)。扫描时只记录 lowest terms 的 fractions，
跳过不是 lowest terms 的重复写法，就得到 \(\Q^+\) 的 enumeration。

最后将
\[
  \Q=\{0\}\cup \Q^+\cup(-\Q^+)
\]
按 \(0,q_1,-q_1,q_2,-q_2,\ldots\) 交错排列。于是 \(\Q\) is countable。
\end{proof}

\begin{explanation}
diagonal argument 在这里解决的是 ``每一行无限长'' 的问题。若逐行扫描，
第一行永远不会结束；按 \(p+q\) 扫描时，第 \(s\) 条 diagonal 只有
\(s-1\) 个 pairs，因此可以在有限时间内进入下一条 diagonal。
\end{explanation}

\subsection{Cardinal inequalities and Cantor--Bernstein}

\begin{definition}[cardinal inequality]
设 \(X,Y\) 是 sets。若存在 injection \(f:X\to Y\)，则写
\[
  \card{X}\le \card{Y}.
\]
若 \(\card{X}\le\card{Y}\) 且 \(\card{X}\ne\card{Y}\)，则写
\(\card{X}<\card{Y}\)。
\end{definition}

\begin{theorem}[Cantor--Bernstein theorem]
若存在 injections \(f:X\to Y\) and \(g:Y\to X\)，则存在 bijection
\(h:X\to Y\)。换言之，
\[
  \card{X}\le\card{Y}
  \quad\text{and}\quad
  \card{Y}\le\card{X}
  \quad\Longrightarrow\quad
  \card{X}=\card{Y}.
\]
\end{theorem}

\begin{explanation}
Cantor--Bernstein theorem 允许我们避免寻找 complicated explicit
bijection。很多时候，两个 injections 很容易构造：例如 inclusion 给出
\((0,1)\to[0,1]\)，linear compression 给出 \([0,1]\to(0,1)\)。theorem
保证这已经足以证明 equal cardinality。
\end{explanation}

\begin{exerciseblock}[Homework 9 Problem 3]
Cantor--Bernstein struggles. Let \(X,Y\) be two sets. Let \(f:X\to Y\) and
\(g:Y\to X\) be injections. Define \(A_0=X\), \(A_{n+1}=g(f(A_n))\), and
\[
  A=\bigcap_{n\ge 0} A_n.
\]
For each proposed function in Homework 9, give an example showing that \(h\)
does not have to be a bijection.
\end{exerciseblock}

\begin{solution}
令 \(X=Y=\N\)，令 \(g=\id_{\N}\)，并定义
\[
  f(0)=0,\qquad f(n)=n+1\quad(n\ge 1).
\]
则 \(f\) is injective，因为 \(0\) 的 image 是 \(0\)，而 positive inputs
\(n\ge 1\) 的 image 是 \(n+1\)，不同 positive inputs 有不同 image。
\(g\) 也是 injection。注意 \(1\notin f(\N)\)，所以 \(f\) is not
surjective。

在这个例子中，
\[
  A_0=\N,\qquad
  A_n=\{0\}\cup\{n+1,n+2,n+3,\ldots\}\quad(n\ge 1),
\]
因此
\[
  A=\bigcap_{n\ge 0}A_n=\{0\},
  \qquad
  A_n\setminus A_{n+1}=\{n+1\}\quad(n\ge 0).
\]

对 Homework 9 Problem 3(a)，proposed map 在 \(A_n\setminus A_{n+1}\)
以及 \(A\) 上都使用 \(f\)。这些 pieces 覆盖整个 \(X=\N\)，所以
\[
  h(n)=f(n)\quad\text{for all }n\in\N.
\]
因此 \(h(\N)=\{0,2,3,4,\ldots\}\)，没有命中 \(1\)。于是 \(h\) is not
surjective，故不是 bijection。

对 Homework 9 Problem 3(b)，proposed map 在 \(A_n\setminus A_{n+1}\)
上使用 \(f\)，在 \(A=\{0\}\) 上使用 \(g^{-1}\)。因为 \(g=\id_{\N}\)，
有 \(g^{-1}(0)=0=f(0)\)。所以此时仍有
\[
  h(0)=0,\qquad h(n)=n+1\quad(n\ge 1).
\]
这个 \(h\) 的 image 仍为 \(\{0,2,3,4,\ldots\}\)，没有命中 \(1\)。因此
\(h\) is not surjective，故不是 bijection。
\end{solution}

\begin{explanation}
Cantor--Bernstein 的正确 proof 会精细地区分哪些 \(x\in X\) 应该用
\(f\)，哪些 \(x\in X\) 应该经由 \(g^{-1}\) 回到 \(Y\)。Homework 9 的两个
错误 formula 把 stable intersection \(A\) 处理错了。本 counterexample
让 \(A=\{0\}\)，并让 \(f\) 缺少一个 image point \(1\)，这样 proposed map
会退化成 \(f\) 或几乎退化成 \(f\)，surjectivity 立即失败。
\end{explanation}

\subsection{Products of countable sets}

\begin{exerciseblock}[Worksheet 8 Exercise 2]
Prove that \(\card{\N\times\N}=\card{\N}\).
\end{exerciseblock}

\begin{solution}
定义 Cantor pairing function
\[
  \pi:\N\times\N\to\N,\qquad
  \pi(m,n)=\frac{(m+n)(m+n+1)}{2}+n.
\]
设 \(s=m+n\)。pair \((m,n)\) 被放在第 \(s\) 条 diagonal 上，而
\[
  T_s=\frac{s(s+1)}{2}
\]
是前面所有 diagonals 的总长度。第 \(s\) 条 diagonal 中，\(n\) 从
\(0\) 到 \(s\) 区分这条 diagonal 上的 \(s+1\) 个 pairs。

若 \(\pi(m,n)=\pi(m',n')\)，令 \(s=m+n\)，\(s'=m'+n'\)。数列
\[
  T_s,T_s+1,\ldots,T_s+s
\]
恰好是第 \(s\) 条 diagonal 的 codes，并且这些 intervals 对不同的
\(s\) 两两不交。因此 \(s=s'\)。在同一条 diagonal 上，code 的最后一项
\(+n\) 决定了 \(n\)，于是 \(n=n'\)，再由 \(m=s-n\) 得 \(m=m'\)。
所以 \(\pi\) is injective。

反过来，给定 \(r\in\N\)，存在唯一的 \(s\in\N\) 使得
\[
  T_s\le r<T_{s+1}.
\]
令 \(n=r-T_s\)，则 \(0\le n\le s\)，令 \(m=s-n\)。于是
\(\pi(m,n)=r\)。所以 \(\pi\) is surjective。综上，\(\pi\) is bijective，
故 \(\card{\N\times\N}=\card{\N}\)。
\end{solution}

\begin{explanation}
这个 proof 是 diagonal enumeration 的 formula 版本。每个 pair
\((m,n)\) 有唯一的 diagonal number \(m+n\)，每条 diagonal 有有限多个
pairs，因此可以把二维 grid 排成一条 sequence。surjectivity 的关键是：
任意 natural number \(r\) 都落在唯一一个 diagonal 的 code interval 中。
\end{explanation}

\begin{exerciseblock}[Homework 9 Problem 4]
Show that \(\card{\N\times\N\times\N}=\card{\N}\).
\end{exerciseblock}

\begin{solution}
由 Worksheet 8 Exercise 2，存在 bijection
\(\pi:\N\times\N\to\N\)。定义
\[
  F:\N\times\N\times\N\to\N,\qquad
  F(a,b,c)=\pi(\pi(a,b),c).
\]
先证明 \(F\) is injective。若
\[
  F(a,b,c)=F(a',b',c'),
\]
则
\[
  \pi(\pi(a,b),c)=\pi(\pi(a',b'),c').
\]
因为 \(\pi\) is injective，有 \(c=c'\) 且
\(\pi(a,b)=\pi(a',b')\)。再次使用 \(\pi\) 的 injectivity，得到
\(a=a'\) and \(b=b'\)。所以 triples 相同，\(F\) is injective。

再证明 \(F\) is surjective。取任意 \(r\in\N\)。因为 \(\pi\) is
surjective，存在 \(u,c\in\N\) 使得 \(r=\pi(u,c)\)。再次使用
surjectivity，存在 \(a,b\in\N\) 使得 \(u=\pi(a,b)\)。于是
\[
  r=\pi(\pi(a,b),c)=F(a,b,c).
\]
故 \(F\) is surjective。综上 \(F\) is bijective，因此
\[
  \card{\N\times\N\times\N}=\card{\N}.
\]
\end{solution}

\begin{explanation}
三重 product 不需要重新发明三维 enumeration。先把前两个 coordinates
\((a,b)\) pair 成一个 natural number，再把这个 number 与第三个
coordinate \(c\) pair。bijection 的 injectivity and surjectivity 在
composition 中逐层传递。
\end{explanation}

\section{Cantor's theorem and uncountability}

\begin{theorem}[Cantor's theorem]
对任意 set \(X\)，
\[
  \card{X}<\card{\Pow(X)}.
\]
\end{theorem}

\begin{proof}
function \(i:X\to\Pow(X)\), \(i(x)=\{x\}\), is injective，所以
\(\card{X}\le\card{\Pow(X)}\)。接着证明不存在 bijection
\(X\to\Pow(X)\)。

假设 \(f:X\to\Pow(X)\) is surjective。定义 diagonal set
\[
  T=\setbuilder{x\in X}{x\notin f(x)}.
\]
因为 \(T\subseteq X\)，所以 \(T\in\Pow(X)\)。由 surjectivity，存在
\(y\in X\) 使得 \(f(y)=T\)。现在检查 \(y\) 是否属于 \(T\)：
\[
  y\in T
  \iff y\notin f(y)
  \iff y\notin T.
\]
一个 statement 不可能与自己的 negation 等价，所以假设导致
contradiction。因此不存在 surjection \(X\to\Pow(X)\)，更不存在
bijection。于是 \(\card{X}<\card{\Pow(X)}\)。
\end{proof}

\begin{explanation}
diagonal set \(T\) 的定义故意在第 \(x\) 个位置与 \(f(x)\) 相反：若
\(x\in f(x)\)，则 \(x\notin T\)；若 \(x\notin f(x)\)，则 \(x\in T\)。
因此 \(T\) 与 list 中每一个 \(f(x)\) 至少在 \(x\) 这个 membership
question 上不同。这个机制是 diagonal argument 的核心。
\end{explanation}

\begin{theorem}[The reals are uncountable]
\[
  \card{\N}<\card{\R}.
\]
\end{theorem}

\begin{proof}
由 Cantor's theorem，
\[
  \card{\N}<\card{\Pow(\N)}.
\]
构造 injection \(\Phi:\Pow(\N)\to[0,1]\subseteq\R\)。对
\(S\subseteq\N\)，令
\[
  \Phi(S)=\sum_{n=0}^{\infty}\frac{a_n}{3^{n+1}},
  \qquad
  a_n=
  \begin{cases}
    1,&n\in S,\\
    0,&n\notin S.
  \end{cases}
\]
若 \(S\ne S'\)，取最小的 \(k\) 使得 \(a_k\ne a'_k\)。前 \(k\) 项相同，
第 \(k\) 项贡献的差的绝对值是 \(3^{-(k+1)}\)。后面所有项最多能抵消
\[
  \sum_{n=k+1}^{\infty}\frac{1}{3^{n+1}}
  =\frac{1}{2\cdot 3^{k+1}},
\]
严格小于第 \(k\) 项差。因此 \(\Phi(S)\ne\Phi(S')\)，\(\Phi\) is
injective。于是 \(\card{\Pow(\N)}\le\card{\R}\)。结合
\(\card{\N}<\card{\Pow(\N)}\)，得到 \(\card{\N}<\card{\R}\)。
\end{proof}

\begin{explanation}
这里使用 ternary digits \(0\) and \(1\)，而不是 binary digits，是为了
避免不同 expansions 产生同一个 real number 的麻烦。第一个不同 digit
造成的差大于后续 tail 的总可能抵消量，所以 encoding 是 injective。
\end{explanation}

\begin{exerciseblock}[Worksheet 8 Exercise 3]
Let \([0,1]=\{x\in\R:0\le x\le 1\}\). Which of the following is true?
\[
  \card{\N}>\card{[0,1]},\qquad
  \card{\N}=\card{[0,1]},\qquad
  \card{\N}<\card{[0,1]}.
\]
\end{exerciseblock}

\begin{solution}
正确关系是
\[
  \card{\N}<\card{[0,1]}.
\]
由 Cantor's theorem 有 \(\card{\N}<\card{\Pow(\N)}\)。上一个 theorem 的
proof 实际构造了 injection \(\Pow(\N)\to[0,1]\)，所以
\[
  \card{\N}<\card{\Pow(\N)}\le\card{[0,1]}.
\]
如果 \(\card{\N}=\card{[0,1]}\)，则由 transitivity 得
\(\card{\Pow(\N)}\le\card{\N}\)，这会与
\(\card{\N}<\card{\Pow(\N)}\) 矛盾。因此 \(\card{\N}\ne\card{[0,1]}\)。
结合 \(\card{\N}\le\card{[0,1]}\)，得到
\(\card{\N}<\card{[0,1]}\)。
\end{solution}

\begin{explanation}
这道题的易错点是把 \([0,1]\) 的有限长度误读成有限或 countable。
cardinality 不测量 length；它只问是否存在 bijection。Cantor's theorem
说明 even the unit interval already has strictly more points than \(\N\)。
\end{explanation}

\section{Continuum hypothesis \optional}

\begin{exerciseblock}[Lecture Notes Section 6.2 Continuum hypothesis]
State the Continuum Hypothesis and explain its role in comparing
\(\card{\N}\), intermediate cardinalities, and \(\card{\R}\).
\end{exerciseblock}

\begin{solution}
Continuum Hypothesis, abbreviated CH, states that there is no set \(S\) such
that
\[
  \card{\N}<\card{S}<\card{\R}.
\]
Equivalently, if \(\aleph_0=\card{\N}\) and
\(\mathfrak c=\card{\R}\), then CH says that the next uncountable cardinal is
already \(\mathfrak c\). In the common notation \(\mathfrak c=2^{\aleph_0}\),
CH is often written as
\[
  2^{\aleph_0}=\aleph_1.
\]

The lecture notes emphasize that CH is not a theorem of the usual ZFC axioms.
If ZFC is consistent, then adding CH does not produce a contradiction, and
adding the negation of CH also does not produce a contradiction. Godel proved
the consistency direction for CH relative to ZFC, and Cohen proved the
corresponding consistency direction for \(\lnot\mathrm{CH}\) using forcing.
Thus, within this course, CH should be treated as an important independent
statement about cardinality, not as a tool that may be freely used in ordinary
proofs.
\end{solution}

\begin{explanation}
Cantor's theorem gives a strict jump \(\card{\N}<\card{\Pow(\N)}\), and the
standard encodings identify \(\card{\Pow(\N)}\) with \(\card{\R}\). CH asks
whether there is any cardinality strictly between those two. The answer cannot
be forced by the usual axioms alone, so a proof in this course should not assume
CH unless the problem explicitly grants it.
\end{explanation}

\section{Axiom of Choice and Zorn's Lemma \optional}

\begin{definition}[choice function]
Let \(\mathcal S\) be a set whose elements are non-empty sets. A choice
function for \(\mathcal S\) is a function
\[
  c:\mathcal S\to\bigcup\mathcal S
\]
such that \(c(A)\in A\) for every \(A\in\mathcal S\).
\end{definition}

\begin{theorem}[Axiom of Choice]
If \(\mathcal S\) is a set whose elements are non-empty sets, then
\(\mathcal S\) has a choice function.
\end{theorem}

\begin{exerciseblock}[Lecture Notes Section 6.3 Axiom of choice and Zorn's Lemma]
Explain the Axiom of Choice, two cardinality consequences, and the statement of
Zorn's Lemma.
\end{exerciseblock}

\begin{solution}
The Axiom of Choice says that for any family \(\mathcal S\) of non-empty sets
one may choose one element from each member simultaneously. Formally, there is a
choice function \(c\) with \(c(A)\in A\) for every \(A\in\mathcal S\). For a
finite family this can be done by finitely many ordinary choices; the axiom is
needed when no explicit rule is available for infinitely many choices.

First consequence: if \(f:X\to Y\) is surjective, then
\(\card{Y}\le\card{X}\). For every \(y\in Y\), the fiber
\[
  f^{-1}(\{y\})=\setbuilder{x\in X}{f(x)=y}
\]
is non-empty. By the Axiom of Choice, choose \(g(y)\in f^{-1}(\{y\})\) for
each \(y\). Then \(g:Y\to X\) is injective: if \(g(y)=g(y')=x\), then
\[
  y=f(x)=y'.
\]
Therefore \(\card{Y}\le\card{X}\).

Second consequence: a countable union of countable sets is countable. Let
\(\{A_n\}_{n\in\N}\) be countable sets. For each \(n\), choose a surjection
\(f_n:\N\to A_n\); for finite non-empty \(A_n\), repeat one element after all
elements have been listed, and if \(A_n=\emptyset\), omit it from the union.
Choice is used to select the whole family \(\{f_n\}\) at once. Define
\[
  F:\N\times\N\to\bigcup_{n\in\N}A_n,\qquad F(n,m)=f_n(m).
\]
This \(F\) is surjective. Since \(\N\times\N\) is countable, there is a
surjection \(\N\to\N\times\N\). Composing gives a surjection from \(\N\) onto
\(\bigcup_n A_n\), so the union is countable.

Zorn's Lemma is the following statement, equivalent to the Axiom of Choice. Let
\((P,\le)\) be a non-empty partially ordered set. If every chain
\(C\subseteq P\) has an upper bound in \(P\), then \(P\) has a maximal element.
Here a chain means a subset in which any two elements are comparable; an upper
bound of \(C\) is an element \(u\in P\) satisfying \(c\le u\) for all
\(c\in C\); a maximal element \(m\) means there is no \(p\in P\) with \(m<p\).
\end{solution}

\begin{explanation}
The Axiom of Choice turns many local non-empty statements into one global
function. In the surjection result, the local objects are fibers
\(f^{-1}(\{y\})\). In the countable-union result, the local objects are
enumerations of the sets \(A_n\). Zorn's Lemma is the order-theoretic form of
the same principle: instead of choosing elements directly, one proves that every
totally ordered partial construction has an upper bound, then obtains a maximal
construction.
\end{explanation}

\begin{reviewbox}{Zorn proof pattern}
A typical Zorn proof has four steps. Choose \(P\), usually a set of partial
objects ordered by inclusion. Verify \(P\ne\emptyset\). Prove that every chain
has an upper bound, often by taking the union of the chain. Apply Zorn's Lemma
to obtain a maximal object. The final step is to prove that this maximal object
has the desired global property; otherwise it could be extended, contradicting
maximality.
\end{reviewbox}

\section{Intervals}

\begin{theorem}
\[
  \card{(0,1)}=\card{\R}.
\]
\end{theorem}

\begin{proof}
The function
\[
  f:(0,1)\to\R,\qquad
  f(x)=\tan\bigl(\pi(x-\tfrac12)\bigr)
\]
is bijective. Indeed, \(x-\tfrac12\) maps \((0,1)\) onto
\((-\tfrac12,\tfrac12)\), multiplication by \(\pi\) maps this interval onto
\((-\pi/2,\pi/2)\), and tangent maps \((-\pi/2,\pi/2)\) bijectively onto
\(\R\). The composition is therefore a bijection.
\end{proof}

\begin{exerciseblock}[Worksheet 8 Exercise 4]
Let \((0,1)=\{x\in\R:0<x<1\}\). Prove that
\(\card{(0,1)}=\card{[0,1]}\).
\end{exerciseblock}

\begin{solution}
Use Cantor--Bernstein theorem. The inclusion map
\[
  i:(0,1)\to[0,1],\qquad i(x)=x
\]
is injective, so \(\card{(0,1)}\le\card{[0,1]}\)。Conversely, define
\[
  j:[0,1]\to(0,1),\qquad j(x)=\frac{x+1}{3}.
\]
For \(0\le x\le 1\), one has \(1/3\le j(x)\le 2/3\), so \(j(x)\in(0,1)\)。
If \(j(x)=j(y)\), then \((x+1)/3=(y+1)/3\)，hence \(x=y\)。Thus \(j\) is
injective, and \(\card{[0,1]}\le\card{(0,1)}\)。
Cantor--Bernstein theorem gives
\[
  \card{(0,1)}=\card{[0,1]}.
\]
\end{solution}

\begin{explanation}
Adding endpoints \(0\) and \(1\) changes topology and compactness, but it does
not change cardinality. The injection \(j(x)=(x+1)/3\) places the whole closed
interval inside the smaller open interval \((1/3,2/3)\)，which is why
Cantor--Bernstein applies.
\end{explanation}

\begin{exerciseblock}[Lecture Notes Exercise 63]
Show that all intervals have the same cardinality.
\end{exerciseblock}

\begin{solution}
The literal sentence needs the standard non-degenerate interpretation. Empty
intervals and singleton intervals do not have the same cardinality as
\((0,1)\)。For non-empty intervals with at least two points, all usual real
intervals have cardinality \(\card{\R}\)。

First, every bounded open interval \((a,b)\) with \(a<b\) is in bijection with
\((0,1)\) by
\[
  \ell:(0,1)\to(a,b),\qquad \ell(x)=a+(b-a)x.
\]
The inverse is \(y\mapsto (y-a)/(b-a)\)，so \(\ell\) is bijective.

Second, a bounded closed or half-closed interval with \(a<b\) has the same
cardinality as \((0,1)\)。For example, \((a,b)\subseteq[a,b]\) gives an
injection \((a,b)\to[a,b]\)。The map
\[
  [a,b]\to(a,b),\qquad x\mapsto a+\frac{b-a}{4}+\frac{x-a}{2}
\]
has image \([a+(b-a)/4,a+3(b-a)/4]\subset(a,b)\) and is injective. By
Cantor--Bernstein, \(\card{[a,b]}=\card{(a,b)}\)。The same compression works
for \([a,b)\) and \((a,b]\)。

Third, unbounded intervals have cardinality \(\card{\R}\)。For instance,
\[
  (a,\infty)\to\R,\qquad x\mapsto \log(x-a)
\]
is bijective, and \((-\infty,b)\to\R\), \(x\mapsto \log(b-x)\), followed by
negation if desired, gives the left-unbounded case. Endpoints can again be
added or removed by the same Cantor--Bernstein compression argument.

Since \((0,1)\) is bijective with \(\R\)，all non-degenerate standard intervals
are mutually equipotent.
\end{solution}

\begin{explanation}
The proof separates intervals by shape: bounded open, bounded with endpoints,
and unbounded. Affine maps handle bounded open intervals. Cantor--Bernstein
handles endpoints because a finite endpoint change can be absorbed by an
injection into a smaller subinterval. Logarithm handles rays because it sends a
positive distance from the endpoint onto all real numbers.
\end{explanation}

\section{Cantor set \optional}

\begin{exerciseblock}[Lecture Notes Section 6.5 Cantor set]
Construct the Cantor set, describe it by ternary expansions, compare its length
intuition with its cardinality, and prove \(\card{C}=\card{\R}\).
\end{exerciseblock}

\begin{solution}
Start with
\[
  C_0=[0,1].
\]
At stage \(1\), remove the open middle third \((1/3,2/3)\), leaving
\[
  C_1=[0,1/3]\cup[2/3,1].
\]
At each later stage, remove the open middle third of every closed interval that
remained at the previous stage. Thus \(C_n\) is a union of \(2^n\) closed
intervals, each of length \(3^{-n}\)。The Cantor set is
\[
  C=\bigcap_{n=0}^{\infty}C_n.
\]

Every \(x\in[0,1]\) has a ternary expansion
\[
  x=\sum_{k=1}^{\infty}\frac{a_k}{3^k},
  \qquad a_k\in\{0,1,2\}.
\]
When a point has two ternary expansions, choose the expansion that avoids an
eventual tail of all \(2\)'s when possible. Under this convention, \(x\in C\)
exactly when \(x\) has a ternary expansion using only digits \(0\) and \(2\)。
The reason is stage-by-stage: the first removed middle third consists of points
whose first ternary digit is \(1\); at stage \(n\), the removed middle thirds
are exactly the points whose first digit equal to \(1\) occurs at position
\(n\), after the previous digits have avoided \(1\)。

The removed total length is
\[
  \frac13+\frac{2}{9}+\frac{4}{27}+\cdots
  =
  \sum_{n=1}^{\infty}\frac{2^{n-1}}{3^n}
  =
  \frac13\sum_{k=0}^{\infty}\left(\frac23\right)^k
  =1.
\]
So length intuition says the remaining set is extremely small.

Now prove its cardinality. Let \(\{0,1\}^{\N_+}\) be the set of infinite binary
sequences indexed by positive integers. Define
\[
  F:\{0,1\}^{\N_+}\to C,\qquad
  F((s_1,s_2,\ldots))
  =
  \sum_{k=1}^{\infty}\frac{2s_k}{3^k}.
\]
This maps binary digit \(0\) to ternary digit \(0\)，and binary digit \(1\) to
ternary digit \(2\)。The image lies in \(C\)。If two binary sequences first
differ at position \(k\)，their ternary expansions first differ by \(2/3^k\)，
while the later tail can contribute at most
\[
  \sum_{j=k+1}^{\infty}\frac{2}{3^j}=\frac{1}{3^k},
\]
which is smaller than \(2/3^k\)。Thus \(F\) is injective. Conversely, every
point of \(C\) has a chosen ternary expansion using only \(0\) and \(2\)，so
dividing each ternary digit by \(2\) gives a binary sequence mapping to that
point. Hence \(F\) is surjective. Therefore
\[
  \card{C}=\card{\{0,1\}^{\N_+}}=\card{\Pow(\N)}=\card{\R}.
\]
\end{solution}

\begin{explanation}
The Cantor set separates two notions of size. Its total removed length is \(1\),
so in measure language it has length \(0\)。But cardinality counts possible
infinite digit sequences. There are as many \(0/2\)-ternary sequences as binary
sequences, and binary sequences are equivalent to subsets of \(\N\)，which have
cardinality \(\card{\Pow(\N)}=\card{\R}\)。
\end{explanation}

\section{Density \optional}

\begin{definition}[dense subset of \(\R\)]
A subset \(S\subseteq\R\) is dense in \(\R\) if for every \(r\in\R\) and every
\(\eps>0\)，there exists \(s\in S\) such that
\[
  |r-s|<\eps.
\]
\end{definition}

\begin{exerciseblock}[Lecture Notes Section 6.6 Density]
Explain density, prove that \(\Z\) is not dense in \(\R\), prove that
\(\Q\) is dense in \(\R\), and define dense in a subset.
\end{exerciseblock}

\begin{solution}
Density asks whether elements of \(S\) can approximate every real number with
arbitrarily small error. It is not a cardinality notion.

The integers are not dense in \(\R\)。Take \(r=1/2\) and \(\eps=1/4\)。For any
\(n\in\Z\)，the number \(n-1/2\) is a half-integer, so
\[
  |n-1/2|\ge 1/2>1/4.
\]
Thus no integer lies within distance \(1/4\) of \(1/2\)，so \(\Z\) is not
dense.

The rationals are dense in \(\R\)。Let \(r\in\R\) and \(\eps>0\)。By the
Archimedean property, choose \(n\in\N\) with \(n>1/\eps\)，so \(1/n<\eps\)。
Choose an integer \(m\) satisfying
\[
  m\le nr<m+1.
\]
Then \(q=m/n\in\Q\)，and
\[
  0\le r-q=r-\frac{m}{n}<\frac1n<\eps.
\]
Therefore \(|r-q|<\eps\)，so \(\Q\) is dense in \(\R\)。

If \(T\subseteq\R\)，then \(S\subseteq T\) is dense in \(T\) if for every
\(t\in T\) and every \(\eps>0\)，there exists \(s\in S\) with
\[
  |t-s|<\eps.
\]
\end{solution}

\begin{explanation}
\(\Q\) is countable but dense; the Cantor set is uncountable but has empty
interior and length \(0\)。These examples show why density cannot be inferred
from cardinality. Density depends on the order and distance inherited from
\(\R\)，not only on how many elements a set has.
\end{explanation}

\begin{exerciseblock}[Lecture Notes Exercise 64]
Prove the Archimedean property for \(\R\): for any \(x\in\R\) with \(x>0\),
there exists \(n\in\N\) such that \(n>x\).
\end{exerciseblock}

\begin{solution}
We prove the statement using the least upper bound property of \(\R\)。Assume
for contradiction that there exists \(x>0\) such that \(n\le x\) for every
\(n\in\N\)。Then \(\N\) is bounded above in \(\R\)。Let
\[
  \alpha=\sup\N.
\]
Because \(\alpha\) is the least upper bound, \(\alpha-1\) is not an upper bound
of \(\N\)。Therefore there exists \(n\in\N\) such that
\[
  n>\alpha-1.
\]
Adding \(1\) gives \(n+1>\alpha\)。Since \(n+1\in\N\)，this contradicts the
fact that \(\alpha\) is an upper bound of \(\N\)。Hence \(\N\) is not bounded
above in \(\R\)。Consequently, for every \(x>0\)，there exists \(n\in\N\) with
\(n>x\)。
\end{solution}

\begin{explanation}
The proof converts Archimedean property into a statement about boundedness. If
some positive \(x\) were at least every natural number, then \(\N\) would have a
supremum \(\alpha\)。The defining property of supremum says anything smaller
than \(\alpha\)，including \(\alpha-1\)，fails to be an upper bound. That gives a
natural number \(n>\alpha-1\)，and then \(n+1\) passes \(\alpha\)，contradicting
the role of \(\alpha\)。
\end{explanation}

\section{Well-ordering \optional}

\begin{definition}[well-ordered set]
Let \((X,\le)\) be a totally ordered set. It is well-ordered if every non-empty
subset \(S\subseteq X\) has a minimum element.
\end{definition}

\begin{exerciseblock}[Lecture Notes Section 6.7 Well-ordering]
Explain well-ordering, prove the basic examples and non-examples, show that
countable sets can be well-ordered, and state the Well-Ordering Theorem.
\end{exerciseblock}

\begin{solution}
A well-ordering is stronger than a total order. A total order compares every
pair of elements. A well-ordering additionally requires every non-empty subset
to have a first element.

Every finite ordinal \(n=\{0,1,\ldots,n-1\}\) is well-ordered by its usual
order. Let \(S\subseteq n\) be non-empty. Since \(S\) is finite, list its
elements as natural numbers below \(n\)。Among these finitely many indices,
there is a smallest one; that element is the minimum of \(S\)。

\((\N,\le)\) is well-ordered. Let \(S\subseteq\N\) be non-empty. Suppose \(S\)
has no minimum. We prove by induction that no \(n\in\N\) belongs to \(S\)。
For \(n=0\)，if \(0\in S\)，then \(0\) would be the minimum of \(S\)。For the
induction step, assume no \(k<n\) lies in \(S\)。If \(n\in S\)，then \(n\)
would be the minimum of \(S\)，because no smaller natural number lies in
\(S\)。This contradicts the supposition. Therefore \(n\notin S\)。By induction,
\(S=\emptyset\)，contrary to the assumption that \(S\) is non-empty. Hence
\(S\) has a minimum.

\((\Z,\le)\) is not well-ordered. The subset \(\Z\) itself has no minimum:
for any \(z\in\Z\)，the integer \(z-1\) also lies in \(\Z\) and satisfies
\(z-1<z\)。

\((\Q^+,\le)\) is not well-ordered. The subset \(\Q^+\) itself has no minimum:
for any \(q\in\Q^+\)，the rational number \(q/2\) also lies in \(\Q^+\) and
satisfies \(q/2<q\)。

Every countable set can be well-ordered. If \(X\) is finite, write
\(X=\{x_0,\ldots,x_{n-1}\}\) and define \(x_i\le_X x_j\) iff \(i\le j\)。
Every non-empty subset has an element with smallest index. If \(X\) is
countably infinite, choose a bijection \(f:\N\to X\)。Define
\[
  x\le_X y
  \quad\text{iff}\quad
  f^{-1}(x)\le f^{-1}(y)\text{ in }\N.
\]
For any non-empty \(S\subseteq X\)，the preimage \(f^{-1}(S)\subseteq\N\) is
non-empty, hence has a minimum \(m\)。Then \(f(m)\) is the minimum of \(S\)
under \(\le_X\)。

The Well-Ordering Theorem states that every set can be well-ordered. It is
equivalent to the Axiom of Choice. One direction is direct: if every set can be
well-ordered and \(\mathcal S\) is a family of non-empty sets, well-order
\(\bigcup\mathcal S\) and choose from each \(A\in\mathcal S\) its minimum
element. The converse direction, from the Axiom of Choice to well-ordering
every set, requires transfinite recursion: repeatedly choose the next element
from the remaining non-empty subset until all elements have been chosen.
\end{solution}

\begin{explanation}
Well-ordering is about the existence of first elements in all non-empty
subsets, not about the familiar visual order of a set. The usual order on
\(\Z\) fails because one can always move left. The usual order on \(\Q^+\)
fails because one can always halve. A countable set can be well-ordered by
transporting the usual well-ordering of \(\N\) through an enumeration.
\end{explanation}

\section{Review checklist}

\begin{itemize}
  \item To prove equal cardinality, construct a bijection or use two injections
  plus Cantor--Bernstein theorem.
  \item To prove countability of a product such as \(\N\times\N\), enumerate
  by diagonals or use Cantor pairing.
  \item To prove uncountability, use Cantor's theorem or a diagonal argument
  that constructs an object missing from any proposed list.
  \item To compare intervals, use affine bijections, compression injections,
  and Cantor--Bernstein theorem.
  \item To use the Axiom of Choice, identify the family of non-empty sets from
  which simultaneous choices are required.
  \item To use Zorn's Lemma, identify the partially ordered set, prove every
  chain has an upper bound, then interpret maximality.
  \item To discuss density or well-ordering, keep the order or metric structure
  separate from cardinality.
\end{itemize}
